\documentclass[acmlarge,11pt,screen,nonacm]{acmart}
\geometry{letterpaper,margin=1in}
\settopmatter{printacmref=false,printccs=false,printfolios=true}
\setcopyright{none}
\renewcommand{\bibfont}{\normalsize}
\newtheorem{theorem}{Theorem}[section]
\newtheorem{lemma}[theorem]{Lemma}
\newtheorem{corollary}[theorem]{Corollary}
\newcommand{\Bool}{\{0,1\}}
\newcommand{\Fp}{\mathbb F_p}
\newcommand{\Q}{\mathbb Q}
\newcommand{\rh}{\rho}
\title{Rectangular Batching for Disjoint Shortest Paths on DAGs}
\author{Zijian Zeng}
\affiliation{\institution{UCSI University}\city{}\country{Malaysia}}
\email{1002266693@ucsiuniversity.edu.my}
\authorsaddresses{\normalsize Author's Contact Information: Zijian Zeng, UCSI University,
\href{mailto:1002266693@ucsiuniversity.edu.my}{1002266693@ucsiuniversity.edu.my}}
\renewcommand{\shortauthors}{Zijian Zeng}
\begin{document}
\begin{abstract}
\normalsize
Given a positive-weight directed acyclic graph and prescribed source--target
pairs, the disjoint shortest paths problem asks for mutually vertex-disjoint
paths, each shortest for its own pair in the original graph. For every fixed
integer $k\geq3$, we give a deterministic detection algorithm with running time
$O_k(m+n+n^{k+0.3})$. Original-weight preprocessing uses $O_k(m+n)$ additions
and comparisons; the remaining bound includes elementary arithmetic on the
new $O_k(\log(n+2))$-bit operands. Thus, in the usual $m=O(n^2)$ input
convention, $k=3$ and $\delta=0.7$ answer the question of Akmal, Vassilevska
Williams and Wein asking for an $O(n^{k+1-\delta})$ algorithm for some fixed
$k\geq3$.

The algorithm accelerates the established minimum-coordinate frontier
recurrence. At each chronological split, a disjoint partition by the target
minimum permits all transitions to be grouped into rectangular Boolean
products. We charge repeated packing, scattering and products over the full
recursion. A fixed rational tensor identity, obtained from an established
rectangular multiplication bound, can be specialized modulo a sufficiently
large prime to give exact Boolean products with bounded operands. An explicit
Strassen-based alternative gives $O_k(m+n+n^{k-2+\log_2 7})$ word operations.
The method also applies to prescribed disjoint paths with label-specific edge
sets sharing a common acyclic order.
\end{abstract}
\maketitle
\clearpage

\section{The result and its context}
In $k$-disjoint shortest paths ($k$-DSP), the input is a graph together with
$k$ ordered terminal pairs $(s_i,t_i)$. The paths must respect the prescribed
pairing, avoid the prohibited vertex intersections, and each have the globally
shortest length for its own pair in the original graph. Minimizing a sum of
path lengths subject to disjointness is a different task: it may allow one
path to become longer to accommodate another. Here shortestness is fixed
before any conflicts are resolved.

For DAGs, the established bound is $O(mn^{k-1})$. Akmal, Vassilevska Williams
and Wein state this bound and ask whether any fixed $k\geq3$ permits a
polynomial saving over the dense $n^{k+1}$ exponent, or whether a standard
hypothesis rules out such a saving \cite{AVW24}, full version, Open Problem~3.
Their two-pair linear-time detection result does not answer that question.
We give a positive answer for every fixed $k\geq3$.

\begin{theorem}\label{thm:main}
For every fixed integer $k\geq3$, prescribed vertex-disjoint shortest-path
detection on a positive-weight DAG with $n$ vertices and $m$ edges is
deterministic in
\[
 O_k(m+n+n^{k+0.3})
\]
time in the original-weight addition/comparison model. New algebraic operands
occupy $O_k(\log(n+2))$ bits. Their elementary arithmetic
cost is included in the displayed bound. This includes strict vertex
disjointness, endpoint-only sharing, and the convention forbidding only
internal/internal intersections, with the relevant convention fixed in advance.
\end{theorem}

For ordinary graphs $m=O(n^2)$, the case $k=3$ is $O(n^{3.3})$, so
$\delta=7/10$ resolves the full existential question. If arbitrary parallel
arcs are listed, the $O(m)$ input scan remains explicit. No bounded-width,
unit-weight, unique-shortest-path, sparse-instance or distances-given
assumption is made. For $B$-bit integer input weights, actual distance
preprocessing costs must be added, for example
$O_k((m+n)(B+\log(n+2)))$ bit operations using elementary additions and
comparisons. Arbitrary input weights are not assumed to fit a word.

\paragraph{What changes in the known algorithm.}
The frontier table has $O_k(n^k)$ entries, but its explicit transitions can
cost $\Theta(n^{k+1})$. We retain that table and the known rule that only its
least coordinate moves. We accelerate the transitions. At a chronological
split with halves $L,R$, every updated target has minimum in $R$. Its moved
coordinate either lies in $R$, or lies later and a context coordinate lies
in $R$. Selecting the first such context coordinate gives exactly $k$
disjoint output classes. Each class becomes a contraction over $L$, with
an entire labelled context kept together in one matrix column. This preserves
path compatibility and makes the total product dimensions small enough.
Charging all dyadic nodes and all repeated data movement is essential.
OpenAI Codex materially assisted the batching construction, proof development
and diagnostic implementation; its full roles are described at the end.

\paragraph{Prior tools and precise distinction.}
Ordered tuple graphs and terminal devices precede this work. Eppstein's
report \cite{Epp95}, Sections~2--4, gives an
$O(mn^{d-1}+K)$ enumeration bound for $K$ feasible $d$-tuples. The
minimum-coordinate recurrence is displayed explicitly in Fleischer, Ge, Li
and Zhu \cite{FGLZ07}, Section~2.1, equations~(1)--(2), where it is attributed
to Perl--Shiloach. B\'erczi and Kobayashi \cite{BK17}, Proposition~10,
use label-specific permitted edges and a tuple graph for shortest paths.
The present proof rederives those semantics to specify exactly what is batched;
it does not claim the recurrence or shortest-path filtering as new.
Chronological closure acceleration using Boolean multiplication is already
central in Valiant \cite{Val75}. The contribution here is the particular
output partition, rectangular context packing, and full global charge for
the labelled frontier problem. Applying generic closure to all $n^k$ states
does not itself establish the claimed bound.

The numerical consequence uses Le Gall and Urrutia's aspect-two exponent
bound \cite{LGU18}, Table~3 and Section~2.2. Later bounds of Vassilevska
Williams, Xu, Xu and Zhou \cite{WXXZ24}, Table~1, are stronger; no current
matrix multiplication record is claimed here. Choudhary, Kumar and Saggi
\cite{CKS26}, Theorem~1.1, instead treat two prescribed pairs in more general
directed graphs. Their $O(mn\log n)$ detection result addresses a different
number of pairs and does not give this fixed-$k\geq3$ DAG bound.

\section{Reduction to labelled paths in a common order}\label{sec:reduction}
Compute distances $d_i$ from each original source by topological relaxation
in the original DAG. If $t_i$ is unreachable, return NO. For label $i$ retain
exactly the edges $(u,v)$ with finite $d_i(u)$ and
\begin{equation}\label{eq:tight}
 d_i(v)=d_i(u)+w(u,v).
\end{equation}
Every shortest $s_i$--$t_i$ path uses only these edges, since each of its
prefixes is shortest. Conversely, summing \eqref{eq:tight} along any retained
$s_i$--$t_i$ path telescopes to $d_i(t_i)$. Thus its original length is
shortest. This costs $O_k(m+n)$ original-weight operations. All retained
graphs share the original topological order; distances are never recomputed
after conflicts or deletions.

First use strict disjointness, including endpoints. Endpoints belonging to
different labels cannot coincide; if they do, return NO. Within one pair,
$s_i=t_i$ is allowed and represents the trivial path. Append a private sink
$z_i$ for each label, after all original vertices, with the sole label-$i$
completion edge $t_i\to z_i$. Appending or removing these edges preserves
feasibility. A path passing through $t_i$ without completing cannot later
return to it in a DAG. All unfinished positions precede all completed sinks.

For the convention allowing endpoint/endpoint sharing only, replace each
terminal vertex $v$ by a private copy $v_i$ for every label having $v$ as an
endpoint. A label without that endpoint cannot visit $v$. Map its retained
edges to its appropriate private copies, forbid entry into its source and
exit from its target except for completion, and leave all other vertices
shared. Endpoint uses separate, while every forbidden internal intersection
remains a collision or a prohibited visit. If only internal/internal
intersections are forbidden, keep in addition one shared internal copy of
each terminal $v$, available precisely to labels for which $v$ is not an
endpoint. Two internal uses then collide; endpoint/internal uses do not.
A trivial pair uses its single endpoint copy and completes.

These transformations are equivalent in both directions by mapping each use
of an original vertex to its designated copy. Only $O_k(1)$ vertices are
added, since there are at most $2k$ terminal vertices and $k$ is fixed.
Copies of a vertex occupy consecutive positions in the common order. An
original edge always joins different original ranks, so every mapped edge
is forward. We have reduced all stated conventions to strict disjointness
in $k$ label-specific forward-edge graphs. Pad this order with isolated
vertices to a power of two $N=O_k(n)$, indexed by $0,\ldots,N-1$.

\section{The frontier recurrence and chronological invariant}\label{sec:frontier}
A valid state is an ordered tuple $q=(q_1,\ldots,q_k)$ of pairwise distinct
positions, with $\rh(q)=\min_i q_i$. Put $F(q)=0$ initially, except for the
source tuple, where $F=1$. In the frontier graph only the label $i$ occupying
the unique minimum $u=q_i$ can move, along an allowed edge $(u,x)$, provided
the resulting tuple remains valid. Since $x>u$, the minimum strictly increases.

\begin{lemma}\label{lem:semantics}
The private-sink tuple is reachable in the frontier graph if and only if
there is a feasible family of prescribed labelled disjoint paths.
\end{lemma}
\begin{proof}
Append each frontier move to the corresponding path prefix. Every abandoned
vertex is smaller than every current coordinate. Initially there are none;
after moving the least coordinate, the newly abandoned vertex is smaller
than its replacement and all unchanged coordinates. Thus a later destination
cannot hit an abandoned vertex. The validity check forbids a hit on any
current coordinate, proving soundness.

For completeness fix a feasible family ending at the private sinks. Move
the least token along its path in that family. If any token is unfinished,
an unfinished token is least, because every original vertex precedes every
private sink. Its next edge exists, and disjointness makes the move legal.
The total number of remaining edges decreases, so the sink tuple is reached.
\end{proof}

Let $A_i$ be label $i$'s Boolean adjacency matrix. For a valid target $p$,
every incoming transition moving label $i$ has exactly the form
\begin{equation}\label{eq:recurrence}
 F\bigl(p[i\!:=\!u]\bigr)\ \wedge\ A_i(u,p_i),
 \qquad u<\rh(p).
\end{equation}
Here $p[i\!:=\!u]$ replaces coordinate $i$ by $u$; invalid states always
have value zero. The target value is the OR of these terms and its initial
seed. Each term uses one complete predecessor tuple, so it never merges
witnesses chosen independently for different paths. At intermediate tuples,
$F$ means frontier-graph reachability; no stronger prefix interpretation
is needed.

Recursively finalize states in increasing order of their minimum. At a node
$I=[l,r]$, split into equal halves $L=[l,h]$ and $R=[h+1,r]$ of size $b$.
Finalize $L$, apply all transitions with source minimum in $L$ and target
minimum in $R$, then finalize $R$. A singleton needs no internal transition.
At entry to a call, transitions from minima below $l$ have already been
applied, and all smaller-minimum states are final. Every transition with
source minimum $u$ and target minimum $v>u$ belongs to the unique node whose
split separates $u$ into $L$ and $v$ into $R$. Induction proves that each
transition is applied once, with final source values, before its target is
finalized. A seed later in the order remains present throughout this process.

\section{The output partition and rectangular batching}\label{sec:batching}
At a cross step define $S=[h+1,N-1]$ and $T=[r+1,N-1]$, so $S=R\cup T$.
Fix the moved label $i$ and order the other $k-1$ labels. Write their positions
as one labelled context $c=(c_1,\ldots,c_{k-1})$. For $u\in L$ let
$D(u,c)=F(c\text{ with }u\text{ inserted at label }i)$. Only contexts in
$S^{k-1}$ are considered. Invalid contexts give zero entries. If $x$ is the
new position of label $i$, partition the outputs as follows:
\begin{align}
 \mathcal C_0 &: \quad x\in R,\quad c\in S^{k-1};\label{eq:classes0}\\
 \mathcal C_j &: \quad x\in T,\quad c_j\in R,\quad
 c_a\in T\ (a<j),\quad c_a\in S\ (a>j),
 \quad 1\leq j\leq k-1.\label{eq:classesj}
\end{align}
These classes are disjoint and cover exactly the outputs with minimum in
$R$. If $x\in R$ use $\mathcal C_0$. Otherwise $x\in T$ and the first
context coordinate in $R$ uniquely determines $j$. Targets with repeated
coordinates will be masked out. In particular the moved coordinate may lie
beyond the current interval $I$; excluding such destinations would be wrong.
For $k=3$ the context domains are respectively $S\times S$, $R\times S$
and $T\times R$, with $x\in R,T,T$ respectively.

For one class, batch at most $b$ destinations and at most $b^{k-1}$ contexts.
With zero padding, multiply the Boolean matrices
\begin{equation}\label{eq:product}
 P[x,u]=A_i(u,x)\quad(b\times b),\qquad
 Q[u,c]=D(u,c)\quad(b\times b^{k-1}).
\end{equation}
For every output $(x,c)$, the product is the OR over $u\in L$ of exactly
the corresponding terms in \eqref{eq:recurrence}. Scan the result, discard
invalid target tuples, and OR it into $F$. Inputs have minimum in $L$;
all writes have minimum in $R$. Therefore no product in this cross step
depends on another product's output. Label, class and batch order do not
affect the result.

\begin{samepage}
\begin{theorem}\label{thm:primitive}
For fixed $k\geq2$, let $\mathcal B_k(b)$ be the cost of deterministic exact
Boolean multiplication of a $b\times b$ matrix by a $b\times b^{k-1}$ matrix,
including input and output arrays. The labelled common-order disjoint-path
problem is solvable in
\begin{equation}\label{eq:total}
 O_k\left(N^k\log(N+2)+
 \sum_{b\in\{1,2,4,\ldots,N/2\}}(N/b)^k\mathcal B_k(b)\right).
\end{equation}
Its space is $O_k(N^k)$ words plus the multiplication primitive's workspace.
\end{theorem}
\end{samepage}
\begin{proof}
Class $0$ has at most $N^{k-1}$ contexts and one destination batch, hence
at most $(N/b)^{k-1}$ products. Class $j>0$ has at most $bN^{k-2}$ contexts
and at most $N/b$ destination batches. Its number of context batches is at
most $(N/b)^{k-2}$: this is an integer after power-of-two padding, so rounding
up cannot increase the upper bound. Each class therefore uses at most
$(N/b)^{k-1}$ products. Empty domains cause no calls. Across $k$ labels and
$k$ classes the bound is $k^2(N/b)^{k-1}$ products per node.

Every product packs $b^2+b^k$ entries and scans, masks and scatters $b^k$
entries. Charge $O_k(b^k)$ every time, including when repacking an operand
for a different destination batch. Context enumeration costs
$O_k(N^{k-1})$ per node and fits within
$O_k((N/b)^{k-1}b^k)=O_k(N^{k-1}b)$. Contexts can be enumerated in streaming
batches; there is no full $N^k$ table scan at any one node. Fixed $k$ gives
constant-time tuple indexing on $O_k(\log(N+2))$-bit words.

There are $N/(2b)$ nodes at half-width $b$. Their multiplication costs sum
to $O_k((N/b)^k\mathcal B_k(b))$, and their other data movement costs sum
to $O_k(N^k)$. Add the $O_k(N^k)$ initialization and sum all levels to obtain
\eqref{eq:total}. The table, adjacency matrices, streamed packed arrays and
recursion stack take $O_k(N^k)$ words; any additional primitive workspace
must be added. Correctness follows from the chronological invariant,
\eqref{eq:product}, and Lemma~\ref{lem:semantics}.
\end{proof}

In particular, if $\mathcal B_k(b)=O(b^\tau)$ with fixed $\tau>k$, the
geometric sum in \eqref{eq:total} is $O_k(N^\tau)$ and absorbs
$N^k\log(N+2)$. This strict inequality will hold for both implementations
below. Additional global bit overhead is charged explicitly.

\section{An explicit square-multiplication implementation}\label{sec:fallback}
Partition the $b^{k-1}$ context columns into $b^{k-2}$ blocks of $b$ columns.
Each block is an integer product of two Boolean square matrices, followed by
testing positivity. Strassen's seven-product identities \cite{Str69} give
\[
 \mathcal B_k(b)=O_k(b^{k-2+\alpha}),\qquad \alpha=\log_2 7.
\]
This includes the $b^k$ array entries since $\alpha>2$. The exponent lies
strictly between $k$ and $k+1$. Theorem~\ref{thm:primitive} gives
$O_k(N^{k-2+\alpha})$ word operations, with $O_k(N^k)$ live words when square
products are streamed. This already supplies a polynomial dense saving.

For completeness, exact arithmetic uses bounded operands. At Strassen depth
$d$, input sums have magnitude at most $2^d$. Leaf products have magnitude
at most $b^2$, and each return level combines at most four signed terms per
entry. Thus all intermediate integers have magnitude $O(b^4)$ and occupy
$O(\log(N+2))$ bits. Positivity, rather than parity, gives the Boolean
answer. Elementary multiplication introduces at most a
$O(\log^2(N+2))$ bit factor; this factor is separate from the displayed
word-operation bound. The next implementation proves Theorem~\ref{thm:main}
including that elementary bit cost.

\section{An exact bounded-operand rectangular primitive}\label{sec:rectangular}
We explain the algebraic dependency because a field-operation exponent
alone does not prove the claimed Boolean time. Write
\[
 \langle a,b,c\rangle=\sum_{\mu=1}^a\sum_{\nu=1}^b\sum_{\xi=1}^c
 X_{\mu\nu}Y_{\nu\xi}Z_{\xi\mu}
\]
for the matrix multiplication tensor. Le Gall and Urrutia
\cite{LGU18}, Section~2.2 and Table~3, establish an exponent at most
$3.251640$ for $\langle q,q^2,q\rangle$. Their polynomial-degeneration
formalism is over a field; specialize it to $\Q$. Choosing an exponent
strictly between that bound and $3.26$, and then a sufficiently large fixed
$q_0$, absorbs the asymptotic constant and gives a rational border-rank
decomposition with $r_0<q_0^{3.26}$ scalar products. No infimum is assumed
to be attained.

\begin{lemma}\label{lem:exactbase}
There exist a fixed integer $q$ and one exact rational bilinear algorithm
of rank $r<q^{3.27}$ multiplying a $q\times q$ matrix by a $q\times q^2$
matrix.
\end{lemma}
\begin{proof}
Fix the preceding polynomial degeneration. Its leading nonzero coefficient
is $\lambda^s\langle q_0,q_0^2,q_0\rangle$ and its degree is at most $D$,
for fixed $s,D$. Its $t$-fold tensor power has leading coefficient at degree
$st$ equal to $\langle q_0^t,q_0^{2t},q_0^t\rangle$ and degree at most $Dt$.
Interpolate this coefficient from $Dt+1$ distinct rational evaluations.
Each evaluation has rank at most $r_0^t$, so exact rational rank is at most
$(Dt+1)r_0^t$. For a sufficiently large fixed $t$, this is less than
$(q_0^t)^{3.27}$. Put $q=q_0^t$.

Tensor rank is invariant under permuting its three variable groups. To make
the required orientation explicit, for indices $a,c\in[q]$, $b\in[q^2]$,
set
\[
 X'_{ac}=Z_{ca},\qquad Y'_{cb}=Y_{bc},\qquad Z'_{ba}=X_{ab}.
\]
Then $\sum_{a,b,c}X_{ab}Y_{bc}Z_{ca}
=\sum_{a,c,b}X'_{ac}Y'_{cb}Z'_{ba}$, which is
$\langle q,q,q^2\rangle$. This exchanges the first and third variable
groups and reindexes them. Reading the coefficients of the third group
gives the required bilinear algorithm with the same exact rational rank.
\end{proof}

All coefficients of this one base identity are fixed rational constants.
Let $H$ be the product of their nonzero denominators, and choose a fixed
integer $C>|H|+2$. For the $N$-vertex instance choose a prime
$p\in(K,2K)$ with $K=\max(N,C)$, using a deterministic sieve and Bertrand's
theorem. No denominator vanishes modulo $p$, so reducing the base identity
modulo $p$ preserves it exactly. One prime is reused throughout the instance.
Its construction costs at most $O(K\log^2(K+2))$ elementary bit operations,
which is within the final bound, with constants depending on the fixed base.
Reducing the fixed coefficients is also a fixed setup cost.

To multiply a $b\times b$ matrix by a $b\times b^2$ matrix, pad to
$M=q^a$ with $b\leq M<qb$ and recurse using the fixed base over $\Fp$.
If necessary add zero products to arrange $r>q^3$ while retaining
$r<q^{3.27}$; the fixed base can be enlarged in advance. All dense linear
maps and output combinations then satisfy
\[
 T(M)\leq rT(M/q)+O_q(M^3),\qquad
 T(M)=O_q(M^{\log_q r})=O_q(b^{3.27}).
\]
The additive term includes every input linear combination and all outputs;
padding costs $O_q(b^3)$ and is included. Residues use
$O(\log(N+2))$ bits, with elementary modular multiplication and reduction
costing $O(\log^2(N+2))$ bit operations. The original Boolean factors have
integer output entries in $[0,b]$, even after zero padding. Since
$p>N\geq b$, a residue is nonzero exactly when its original integer entry
is positive. Hence testing nonzero computes the Boolean product exactly,
without parity cancellation or random error. We have proved, for this
global modulus,
\begin{equation}\label{eq:rectcost}
 \mathcal B_3(b)=O(b^{3.27}\log^2(N+2)).
\end{equation}
The dependence on $N$ in this implementation is solely the global operand
bit bound; a standalone primitive could choose its own prime above $b$.

For $k>3$, divide the $b^{k-1}$ columns into $b^{k-3}$ blocks of $b^2$
columns and use \eqref{eq:rectcost}. The resulting cost is
\[
O_k(b^{k+0.27}\log^2(N+2)),
\]
including the $b^k$ array entries. Substitution
in \eqref{eq:total} yields
\[
 O_k\bigl(N^k\log(N+2)+N^{k+0.27}\log^2(N+2)\bigr)
 =O_k(N^{k+0.3}).
\]
The prime setup is smaller. Combining this with Section~\ref{sec:reduction}
and $N=O_k(n)$ proves Theorem~\ref{thm:main}.

The rational base may be extremely large. It is one finite identity,
hardwired independently of graph size, so the construction defines a uniform
algorithm rather than size-dependent advice. No coefficient list or fast
rectangular implementation is supplied here. The exponent depends on the
cited degeneration theorem and the exact conversion just proved. The
explicit algorithm of Section~\ref{sec:fallback} has no such implementation
dependency. Interpolation, tensor permutation and prime specialization are
established algebraic principles; no improvement to the underlying matrix
multiplication exponent is asserted.

\section{Verification, scope and limitations}\label{sec:limits}
The proof applies to any fixed collection of label-specific forward-edge
graphs on a common order, not only tight-edge graphs. In particular it also
applies to ordinary prescribed disjoint paths on DAGs. It preserves all
labels and detects feasibility; no search-to-decision or reporting bound
is asserted. The $N^k$ frontier table remains, and dense-instance speedup
does not imply an improvement over every sparse-instance baseline. No
extension to cyclic digraphs, optimal exponent or matching lower bound is
proved.

Exact finite checks accompany the argument. The research implementation
compared 145 complete frontier reached sets on arbitrary labelled and weighted
instances, including $k=2,3,4,5$. A separately written implementation compared
4,718 reached sets: all triples of labelled forward-edge graphs on three
vertices with all ordered distinct seeds (3,072 cases), all common-edge DAGs
on four vertices with ordered three-token seeds (1,536 cases), and 110 further
cases. It also compared 120 weighted original instances under three endpoint
conventions, giving 360 original shortest-path comparisons; cases included
181-bit weights, ties, unreachable targets, repeated terminals and trivial
pairs. All comparisons matched. These finite tests check indexing and
semantics. They do not establish the asymptotic bound, implement the fast
tensor primitive, or replace any proof. The executable checks and their
precise finite coverage are retained with the research artifacts.

\paragraph{AI methodology.}
OpenAI Codex materially assisted the research program, including the
batching construction, proof development, counterexample searches,
verification code, primary-literature comparison, and drafting of all
sections. Separate Codex contexts performed an internal mathematical review
and wrote an independent finite verifier. These model-assisted checks are
not external peer review or a formal proof certificate. The manuscript's
mathematical justification is the argument and identified dependencies above;
no claim of independent human verification is made here. The listed human
author retains responsibility for any use or dissemination of this work.

\clearpage
\bibliographystyle{ACM-Reference-Format}
\bibliography{references}
\end{document}
